\documentclass[12pt]{article}
\usepackage[utf8]{inputenc}
\usepackage[T1]{fontenc}
\usepackage{hyperref}
\title{Tagalaka Language Resources for K-12 Curriculum Development}
\author{Compiled by Claude Sonnet 5 (claude-sonnet-5) for Andr\'as Kornai}
\date{}
\begin{document}
\maketitle
\section{Introduction}
Tagalaka (also Takalak; ISO 639-3: tgz; AIATSIS Y125) is a poorly attested, extinct Pama-Nyungan language of Queensland, traditionally associated with the middle Gilbert River and the Georgetown--Forsayth--Croydon area. A Tagalaka Aboriginal Corporation exists as a community body.
\section{Primary Educational Resources}
No primers, textbooks or school programmes were found.
\section{Linguistic and Reference Materials}
\begin{itemize}
\item Wikipedia, Tagalaka language: \url{https://en.wikipedia.org/wiki/Tagalaka_language}
\item Wikipedia, Tagalag people: \url{https://en.wikipedia.org/wiki/Tagalag_people}
\item G. Breen (1981), The Mayi Languages of the Queensland Gulf Country (catalogue record). Whether it contains Tagalaka material was not confirmed from the search results: \url{https://catalog.hathitrust.org/Record/000106521/Cite}
\item PHOIBLE source entry for Breen 1981: \url{https://phoible.org/sources/breen_mayi_1981}
\item University of Sydney Archives, Tagalaka language record (no vocabulary content seen): \url{https://archives-search.sydney.edu.au/nodes/view/18684}
\end{itemize}
\section{Religious and Community Resources}
The Tagalaka Aboriginal Corporation (tagalaka.com/who-we-are) is cited by Wikipedia and About the North; its content was not retrieved. About the North page: \url{https://aboutthenorth.au/styled-59/styled-184/styled-238/styled-362/}
\section{Digital and Audio Resources}
None found.
\section{Curriculum Development Context}
Documentation is very thin. Any curriculum should be led by the Tagalaka community and begin with an audit of archival holdings (for example AIATSIS and Breen's material) to see what is actually recoverable.
\end{document}